% Generated by Listas/make_statements.py from pratica-lecture-08.md, item 10. Do not edit.
\textbf{Source.} Tadelis (2013), \emph{Game Theory: An Introduction}, Exercise 15.5 ``Entry (again)'', pp.~316--317 {[}PDF pp.~335--336{]}. Reproduced verbatim. Answer: none in the book.



\begin{quote}

\textbf{Entry (again):} The market for widgets has an incumbent firm. The total value of having 100\% of the market is 10, which the incumbent receives if no one enters. A potential entrant arrives, and it can be one of two types: tough (\(T\)) or weak (\(W\)). A weak entrant can choose one of three options: small entry (\(S\)), big entry (\(B\)), or exit (\(X\)). A tough entrant can choose only between \(S\) or \(B\) (it is inconceivable that it would choose \(X\)). There is no cost for a tough entrant to enter at any level. However, it costs a weak entrant 6 to enter at any level. Exiting costs nothing to the entrant. The entrant knows his type, but the incumbent knows only the prior distribution: \(\Pr\{T\} = \frac{1}{2}\).



In response to any level of entry, the incumbent can choose to accommodate (\(A\)) or fight (\(F\)). Accommodating an entrant imposes no costs. Independent of the entrant's type, accommodating small entry gives the incumbent 60\% of the market and the entrant 40\%, while accommodating big entry gives the incumbent 40\% of the market and the entrant 60\%. Fighting a tough entrant increases the incumbent's market share by 20\% (relative to accommodating) but imposes a cost of 4 on the incumbent. Fighting a weak entrant that chose \(S\) increases the market share of the incumbent to 100\% but imposes a cost of 2 on the incumbent. Fighting a weak entrant that chose \(B\) increases the market share of the incumbent to 100\% but imposes a cost of 8 on the incumbent.



\begin{enumerate}

\def\labelenumi{\alph{enumi}.}

\item

  Draw this game in extensive form.

\item

  Using a matrix representation, find all the pure-strategy Bayesian Nash equilibria for this game.

\item

  Which one of the Bayesian Nash equilibria is preferred by the incumbent? Can it be supported as a perfect Bayesian equilibrium?

\item

  Find all the perfect Bayesian equilibria of this game.

\end{enumerate}

\end{quote}

